\scanpage{30}

\deleted{qui décrivent $\unsure{\mathfrak S_3|\mathcal S_3}\times\mathbb Z/2$, on va relever ces générateurs en des éléments de $\unsure{\Sigma|E}$, de façon que $l_\infty,l_1,l_0$ s'expriment comme des mots en ces éléments, et il suffira alors d'ajouter cette relation $l_\infty l_1l_0$. aux relations correspondant aux précédentes \ldots}

$\dot\sigma_0,\dot\sigma_1,\dot\sigma_\infty,\deleted{\unsure{\rho}}\dot\rho,\tau$, avec les relations

\[
\begin{aligned}\bullet\\\bullet\\\bullet\\\bullet\\\bullet\end{aligned}\,
\left\{\begin{aligned}
\dot\sigma_0^2&=\dot\sigma_1^2=\dot\sigma_\infty^2=1\\
\dot\sigma_0\dot\sigma_1&=\dot\sigma_1\dot\sigma_\infty=\dot\sigma_\infty\dot\sigma_0=\dot\rho\\
\dot\rho^3&=1\\
\tau^2&=1\\
\tau&\ \text{commute à }\dot\sigma_0,\dot\sigma_1,\dot\sigma_3\quad\text{(donc aussi à }\dot\rho\text{)}
\end{aligned}\right.
\]

on les relèvera en $\deleted{\tau_0,\tau_1,\tau_\infty}\ \sigma_0,\sigma_1,\sigma_\infty,\rho,\underbrace{\tau_0,\tau_1,\tau_\infty}_{\text{sur }\tau}$ avec les relations

% slanted
\margin{\deleted{NB} Il faut ajouter les relations $\left\{\begin{aligned}\rho\sigma_0\rho^{-1}&=\sigma_1\\\rho\sigma_1\rho^{-1}&=\sigma_\infty\\\rho\sigma_\infty\rho^{-1}&=\sigma_0\end{aligned}\right.$ \deleted{\unsure{sont}} conséquence\unsure{s} des relations ci-contre}

\margin{[ NB \unsure{On posera|On trouve} $l_0=\tau_\infty\tau_1$, $\deleted{l_1=\tau_0\tau_\infty}$, $l_1=\tau_0\tau_\infty$ $\deleted{\unsure{l_\infty={}}}$ $l_\infty=\tau_1\tau_0$ et $l_\infty l_1l_0=1$ automatique ]\\
$\uparrow$ \deleted{\unsure{les}} Ceci implique que $\tau_0,\tau_1,\tau_\infty$ engendrent \unsure{un ss-groupe engendré|le groupe contenant} par $l_0,l_1,l_\infty$}

\[
\left\{\begin{aligned}
&\sigma_0^2=\sigma_1^2=\sigma_\infty^2=\tau_0^2=[\tau_1^2=\tau_\infty^2]=1\\
&\rho^3=1\\
&\rho\tau_0\rho^{-1}=\tau_1,\quad\rho\tau_1\rho^{-1}=\tau_\infty,\quad\rho\tau_\infty\rho^{-1}=\tau_0\\
&\tau_0\sigma_{\deleted{\unsure{\infty}}0}=\sigma_0\tau_0,\quad\tau_1\sigma_1=\sigma_1\tau_1,\quad\tau_\infty\sigma_\infty=\sigma_\infty\tau_\infty\\
&\sigma_0\sigma_1=\tau_\infty\tau_1\rho,\quad\sigma_1\sigma_\infty=\tau_0\tau_\infty\rho,\quad\sigma_\infty\sigma_0=\tau_1\tau_0\rho
\end{aligned}\right.
\]

\margin{Ce sont les éléments $\widetilde\sigma_0,\widetilde\sigma_1,\widetilde\sigma_\infty$ $\longrightarrow$}

\deleted{\uline{Question} \deleted{\unsure{Peut-}} \inserted{On peut} exprimer $\rho$ en fonction des $\sigma_0,\sigma_1,\sigma_\infty,\tau_0,\unsure{1|\tau_1},\tau_\infty$\unsure{??}, \deleted{et non plus} \inserted{par l'une du} dernier groupe de relations, et substituer dans les autres.}
% archive sheet number: 29
